\chapter{CONCLUSION AND FUTURE WORK}
\label{chap:conclusion}

\section{Conclusion}
\label{sec:conclusion}
This project set out to replace the fixed, global number of Householder factors in DeltaProduct with a number chosen per token by a learned mechanism, and to establish how many factors each token requires. The base method was reproduced on a single 6~GB GPU: the $S_3$ results match the published pattern exactly, and the $S_5$ model is exact on 99.8\% of sequences at its training length and reaches a token accuracy of 0.864 at four times that length.

The analysis established that a product of $K$ factors can move at most $K$ dimensions, that factors of fixed reflection strength can only be removed in pairs, and, centrally, that the requirement of a token is set by the cheapest faithful representation of the group generated by the whole alphabet rather than by the token. A 5-cycle requires four factors when the closure is $S_5$ and two when it is $A_5$, and alphabets generating $A_5$ are indeed learned with two factors. With a reference allocation derived from this analysis, per-token order matches the four-factor model while requiring 2.3 to 3.5 times fewer factors per token. A monotone halting gate trained end to end without supervision on the allocation recovered the demand structure exactly in 7 of 49 runs overall, never under-allocated a token and never reduced accuracy; its best runs require 2.20 factors per token at a token accuracy of 0.9998, and one configuration reached the exact allocation of 1.30.

Two findings define the current state of the work. Training in this setting is bistable, so reliability, measured as the fraction of runs that succeed, is now the main limitation of the mechanism. And the reported savings are in required factors per token: the current implementation still executes every factor, so no reduction in time has yet been realised.

\section{Future Work}
\label{sec:future}
\begin{enumerate}[label={(\arabic*)}]
	\item \textbf{Reliability of the learned gate.} Compare the earliest gate configuration, which reached the exact allocation, with later ones one change at a time to find what raised the floor to two factors, and measure the success fraction over enough repeated runs to give narrow confidence intervals.
	\item \textbf{Realised efficiency.} Execute only the allocated factors, for example with a per-level capacity router in the spirit of Mixture-of-Depths \cite{mod} that keeps tensor shapes static, and measure training and inference time against fixed orders at equal accuracy.
	\item \textbf{Multi-level requirements.} Evaluate the gate on alphabets whose requirement takes every value from 1 to 4, and test whether it learns that a 5-cycle requires two factors under an $A_5$ alphabet but four under an $S_5$ one.
	\item \textbf{The high-requirement regime.} Test the removal of the additive input pathway \cite{howe} on the configuration with 96\% maximum-requirement tokens, which remains unresolved.
	\item \textbf{Natural-language state tracking.} Once training is reliable, evaluate the layer on benchmarks such as SCONE \cite{scone} and entity tracking \cite{entitytracking}, where the number of state components an instruction changes gives an independent check on the learned allocation.
\end{enumerate}
